\section{The documents, and what each contains}\label{sec:corpus}

\subsection{The editions and how they relate}

\begin{center}\small
\begin{tabularx}{\linewidth}{@{}>{\raggedright\arraybackslash}p{0.19\linewidth}>{\raggedright\arraybackslash}p{0.23\linewidth}>{\raggedright\arraybackslash}X@{}}
\toprule
Edition & Where & Contents and relation to the others\\
\midrule
Project reader, 488 pages, 7 August 2026 & repository root, with source and Lean archives & The earlier edition. The README says it is preserved and is not the newest reader.\\
Cumulative archive, 619 pages; title page dated 30 August 2026 & Zenodo 22666493 and 22678971, file 00 & The 488-page reader continued. Its main source has the same section skeleton, now 60{,}561 lines instead of 57{,}452. Fourteen satellite files add §§10--19, the $Q=83$ step of §24 (Theorem~24.241) and the carrier refinements §§24.1--24.3. It begins from the source of the earlier Zenodo record 21845035, version~69 (archive p.~3). It is the front file of the Zenodo collection.\\
Bounded congruence transport, 67 pages, 8 September & Zenodo 22666493, files 02--03 & The 58-statement supplement.\\
Expanded supplement, 86 pages, 9 September & Zenodo 22678971, files 07--08; repository \texttt{research/\allowbreak expanded-2026-09-08/} & The same, with 14 further statements: 72 statements in 15 source files, with checkers (§\ref{ssec:supplement}).\\
ES/Fable reader, 124 pages, 20 September & repository \texttt{research/\allowbreak incoming/\allowbreak es-fable-zeta-bridge-\allowbreak 20260920/} & A coordinate bridge from Erdős--Straus solutions to the Fable cover and a weighted conductor of the zeta programme (§\ref{ssec:crosswalk}).\\
Continuation, 17--23 September & repository README, \emph{Start reading} (28 items); \texttt{research/\allowbreak incoming/} & 41 dated packages; the results bench R01--R10 and the integration arguments X1--X7 (§\ref{ssec:continuation}).\\
\bottomrule
\end{tabularx}
\end{center}

These editions have different contents; the repository preserves them rather than overwriting them (README, edition table).

\subsection{The 619-page archive, section by section}\label{ssec:archive}

Page numbers are those of the PDF. The column \emph{Kind} says how a section relates to the equation:
\begin{itemize}
\item \emph{direct}: it proves or tests solvability;
\item \emph{literature}: it audits a published approach to the equation;
\item \emph{structural}: a construction whose connection to solvability is, in the archive's words (p.~4), an interface still to be supplied by an explicit typed map. Section~\ref{sec:overlaps} describes these connections.
\end{itemize}
The status column refers to this paper.

\begin{center}\small
\begin{longtable}{@{}>{\raggedright\arraybackslash}p{0.06\linewidth}>{\raggedright\arraybackslash}p{0.08\linewidth}>{\raggedright\arraybackslash}p{0.45\linewidth}>{\raggedright\arraybackslash}p{0.13\linewidth}>{\raggedright\arraybackslash}p{0.17\linewidth}@{}}
\toprule
§ & pp. & Content & Kind & Status here\\
\midrule
\endhead
front & 3--5 & Research status: the census figures of §24, the four families, the record-shell results, and an explicit list of non-claims & summary & read\\
1 & 6--16 & Scope, road map, mathematical status, attribution, lineage & summary & read in part\\
2 & 16--18 & The discriminant-five identity $(2n+1)^2-5=4(n^2+n-1)$ and its root involution, from u/CivQ17 & structural & located\\
3 & 18--20 & The shell $p_*=8{,}803{,}369$, $R_*=107$, reached through Busy-Beaver values ($S(4)=107$, $\mathrm{Space}(5)=12289$) & direct (one shell) & $R=107$ is the least occupied shell at $p_*$ (Proposition~\ref{prop:records})\\
4 & 20--22 & Alpöge's constant-Jacobian map $F:\C^3\to\C^3$ with a three-point fibre & structural & see \S\ref{sec:overlaps}\\
5 & 22 & Withdrawn constructions (a cube-root interpolation on the Fable fibre), with the point of failure & structural & read\\
6 & 22--27 & Central-angle and inversive-circle theorems of u/UmbrellaCorp\_HR; a Lorentz reflection & structural & located\\
7 & 27--31 & The complete fibre of $F$ over $(-\tfrac14,0,0)$ and a marked-factor lift & structural & see \S\ref{sec:overlaps}\\
8 & 31--38 & The singular $C_1$--$C_2$--$C_3$ support object; dihedral actions & structural & located\\
9 & 38--93 & The exact divisor-shell bijection (Theorem~9.1, for every odd prime) and its layer count (Theorem~9.64, p.~90); then directed blade transport and a quarter-radius inversion & direct (Theorems 9.1 and 9.64), structural (the rest) & the bijection and the count re-proved here (Theorem~\ref{thm:shell}(b))\\
10 & 93--98 & Residual sieve modulo $840$: an occupied shell with $R=3$ or $R=7$ at every prime $p\equiv1\pmod{24}$ outside the six hard classes (Theorem~10.6); the Jacobi-symbol barrier (Theorem~10.9) & direct & proved here (Propositions~\ref{prop:survivors}, \ref{prop:jacobi})\\
11 & 98--100 & Elsholtz--Tao Type~I and Type~II coordinates mapped to the residual shell & literature & located\\
12 & 100--106 & Adjacent programmes (López, Dyachenko, Bradford): what each proves, and where a claimed bridge fails & literature & located\\
13 & 106--114 & Mordell's constructions and Rosati's theorem; Salez's seven charts; Bright--Loughran surfaces; the Vaughan boundary & literature & located\\
14 & 114--119 & López's structural solutions: exact divisor coordinates and a quantifier repair & literature & located\\
15 & 119--124 & The Star--Kneser programme: a one-way forcing theorem and a retracted converse & direct & located\\
16 & 124--127 & The mod-$107$ deficit progression, following u/CommonCareful3149 & direct & progression and its corollary reproduced (§\ref{ssec:records})\\
17 & 127--135 & Two-fraction criteria; Monks--Velingker; Ionascu--Wilson classes, with the printed record list recovered and one printed class label corrected (Proposition~17.6); Zelator's prime-square classification & literature & records recomputed (§\ref{ssec:records})\\
18 & 135--148 & The residual-eleven support ideal, its colons, and the paired $R=7/R=11$ affine theorem & direct & located\\
19 & 148--157 & Strict record shells: an open conjecture and the boundary example $p=806{,}521$ & direct & records and deficits reproduced (§\ref{ssec:records})\\
20 & 157--179 & The shell $p=12289$, $R=11$, $a=3075$: counts $(3,6,3)$, its dihedral orbit, and the Star--Kneser quotient certificate & direct (one shell) & counts reproduced (below the table); $a$ is the least occupied shell; the rest located\\
21 & 179--200 & The split-zero semiring $G(R)=R\sqcup\{\tau\}$ and the $C_1$ carrier & structural & see \S\ref{sec:overlaps}\\
22 & 200--201 & Rigidified denominator towers over $\Ff_1$ & structural & located\\
23 & 201--250 & Mirror completion of the Boolean carrier; the $12289$ orbit matrix & structural & located\\
24 (i) & 250--298 & The opening of §24: Tomita--Takesaki transport and $C_3$ modular data (Theorems 24.1--24.44) & structural & located\\
24 (ii) & 299--604 & The rest of §24, from Theorem~24.45: shell arithmetic.
\begin{itemize}
\item Star--Kneser forcing (Theorems 24.45--24.48, pp.~299--301);
\item endpoint divisor fibres (Theorems 24.50--24.63, pp.~303--312);
\item coupled and adjacent shells, and the base of the census (Theorem~24.100, p.~335);
\item the carriers for $R=11$ to $R=31$, including the constant-factor carriers (Lemma~24.218, p.~422) and the four families (Theorem~24.219, p.~423);
\item the $R=71$, $R=191$ and $Q=83$ steps (Theorems 24.233, 24.237 and 24.241, pp.~436--444);
\item §§24.1--24.3 (pp.~451--477): the $R=167$ and $R=239$ carrier theorems and the refresh of the $R=167$ coefficient;
\item §24.4 (pp.~478--604): a return to the $R=23$ lineage, interleaved with further structural constructions.
\end{itemize}
& direct (with structural material in §24.4) & four families proved here (Theorem~\ref{thm:families}); census located (§\ref{ssec:census})\\
25 & 605--616 & A map from Busy-Beaver extrema to finite shell carriers & structural & see \S\ref{sec:overlaps}\\
\bottomrule
\end{longtable}
\end{center}

At $p=12289$ the shell $a=3075$ has
\[
E_a=\{41,1845,15375\},\qquad M_a=\{5,225,1875,5043,42025,1891125\}.
\]
So $(|E_a|,|M_a|,|E_a|)=(3,6,3)$, and exactly six unordered solutions have least denominator $3075$ (recomputed here).

So about 250 of the 619 pages treat solvability directly:
\begin{itemize}
\item §§3, 9 (Theorems 9.1 and 9.64), 10, 15, 16 and 18--20, about 70 pages;
\item §24 from p.~299 to the end of §24.3 (p.~477), 179 pages.
\end{itemize}
Parts of the 127 pages of §24.4 also do. The literature audits of §§11--14 and 17 add about 35 pages. The rest are structural constructions, whose connection to the equation the archive describes as an interface still to be supplied by an explicit typed map (p.~4).

\subsection{The bounded-transport supplement (72 statements)}\label{ssec:supplement}

The supplement starts from the primes $p=12h+1$ and all shells $3h+1\le a\le9h$ of Theorem~\ref{thm:shell}. Its sections are:
\begin{enumerate}[label=\arabic*.]
\item original shells, coordinates and integer reconstruction, with the first-two-shell sieve and an all-shell no-hit sieve;
\item local congruences with all integral lifts retained, with an integral cyclic bridge;
\item bounded divisibility and simultaneous congruences: unary Boolean transport, a shared-variable Chinese-remainder construction credited to a colleague, raw-scale hit incidence, and a boundary-swap completion;
\item exact cross-shell transport at a fixed rational ratio;
\item a primitive Euclidean successor with its full arithmetic domain;
\item arithmetic localisation along the original norm subdivision, with a reading of Connes' primitive intersections;
\item a strictly decreasing denominator with exact arithmetic return;
\item the Navier--Stokes auxiliary cover: the torus matrix $J=\begin{pmatrix}3&1\\1&5\end{pmatrix}$ of the announced Navier--Stokes manuscript, its covering fibres and character corrections, and sampling and reconstruction estimates.
\end{enumerate}
The supplement's own README states that its results prove or disprove neither the Erdős--Straus conjecture nor the Riemann hypothesis nor the Yang--Mills mass gap, and that they do not validate the Navier--Stokes construction.

\status{Reproduced in part. A second-pass audit, which stopped before it wrote its report, recomputed the finite content of the statements of Sections 1--7. Its scripts, which are independent of the workbench's code, found every checked statement correct:
\begin{itemize}
\item the gluing and Chinese-remainder theorems on random systems and on all shells of small primes;
\item the ratio-transport domain on all $57{,}882$ ordered shell pairs of the primes $p\equiv1\pmod{12}$ below $260$;
\item the shear and no-two-shears theorems on all shells of $131$ primes below $4000$;
\item the Connes marked successor on all primes $p\equiv1\pmod{12}$ below $3\cdot10^5$;
\item the fixed-$y$ successor on all $3{,}114$ source witnesses below $1500$;
\item the orchard partition on all $42$ admissible triples below $1300$;
\item the boundary four-candidate law on all $72{,}282$ raw records below $400$;
\item the unary box, the input/output incidence and the raw-scale theorems on their stated ranges.
\end{itemize}
One of its checks, on the smallest member of a boundary orbit, reported mismatches. The script expected the pattern $(\min(A,B),\max(A,B))$ in every case. The supplement's own table (\texttt{boundary\_swap\_completion.tex}, lines 296--309) gives a one-point orbit's own point as its smallest member, which is what the script found. The mismatch is therefore in the script, not in the supplement. The same pass also looked at scope, and its search led to Proposition~\ref{prop:twoshears} below. The general proofs were not re-derived, except that of Section~5's no-two-shears theorem; the auxiliary-torus Section 8 was not checked. The scripts are in \note{checks/es\_reader/second\_pass\_partial/}.}

Section 5 of the supplement studies the two unimodular shears $L(m,n)=(m+n,n)$ and $J(m,n)=(m,m+n)$ as maps between adjacent shells. It encodes the solutions with first denominator $a$ by the set
\[
 \mathcal W_a=\{(m,n)\in\ZZ_{>0}^2:\ m\mid pa,\ n\mid pa,\ \gcd(m,n)=1,\ R_a\mid m+n\},
\]
which is in bijection with the ordered solutions $(a,y,z)$ through $y=pak/n$, $z=pak/m$, $k=(m+n)/R_a$ (supplement §1, first lemma; this is Theorem~\ref{thm:shell}(b) in other coordinates). A \emph{shear edge} is a step $(a,m,n)\mapsto(a+1,L(m,n))$ or $(a,m,n)\mapsto(a+1,J(m,n))$ with $(m,n)\in\mathcal W_a$ and the image in $\mathcal W_{a+1}$. The supplement proves, for $p=12h+1$ and shells $3h+1\le a\le9h$, that no two shear edges compose (Theorem \emph{No two consecutive positive Euclidean shell steps}, \texttt{publication.tex}, line~933). Its proof uses the hypothesis on $p$ at one step only, and there only through $p\equiv1\pmod3$.

\begin{proposition}[the no-two-shears theorem, exact scope]\label{prop:twoshears}
Let $p$ be a prime, and consider all shells $a$ with $p/4<a<p$.
\begin{enumerate}[label=(\alph*)]
\item If $p\not\equiv2\pmod3$, no two shear edges compose: there is no path $(a,m,n)\mapsto(a+1,m',n')\mapsto(a+2,m'',n'')$ of shear edges.
\item Such paths exist at some primes $p\equiv2\pmod3$. At $p=131$ one is $(33,1,33)\mapsto(34,1,34)\mapsto(35,1,35)$, with $R_a=1,5,9$. At $p=1613$, where $p\equiv1\pmod4$, one is $(404,1,404)\mapsto(405,1,405)\mapsto(406,1,406)$, with $R_a=3,7,11$. Proposition~\ref{prop:sheargraph} says exactly at which primes.
\end{enumerate}
\end{proposition}
\status{\textsf{Generalised here}, with the workbench's proof. The supplement states (a) for $p\equiv1\pmod{12}$ and $3h+1\le a\le9h$ only. Checked on the full graph of every prime below $2000$ (script \note{esr\_twoshears.py}: no path for $p\equiv1\pmod3$; $14$ paths for $p\equiv2\pmod3$), and by a complete search up to $p\le10^7$ that uses the shape of an edge proved below: it finds $19{,}558$ paths of two $L$-edges, all at primes $p\equiv2\pmod3$, and as many of two $J$-edges.}
\begin{proof}
\emph{The shape of one edge} (the supplement's classification, whose proof needs only $a+1<p$). Let $(a,m,n)\mapsto(a+1,m+n,n)$ be an $L$-edge. Then $n$ divides $pa$ and $p(a+1)$, whose greatest common divisor is $p$. If $n=p$, then $p\nmid m$ by coprimality, and $m+p$ divides $p(a+1)$ and is prime to $p$, so $m+p\mid a+1$; but $m+p>p>a+1$. So $n=1$. If $p\mid m$, then $m+1$ is prime to $p$ and exceeds $a+1$, and yet divides $p(a+1)$, which is impossible; so $p\nmid m$ and $m\mid a$. Then $m+1\le a+1<p$, so $m+1\mid a+1$. Finally $R_a\mid m+1$ and $R_{a+1}=R_a+4\mid m+2$. For a $J$-edge, exchange $m$ and $n$; the exchange preserves every $\mathcal W_a$.

(a) An $L$-edge ends at $(a+1,m+1,1)$, whose first coordinate is at least $2$; a $J$-edge starts at a pair whose first coordinate is $1$. So an $L$-edge is never followed by a $J$-edge, and by the exchange a $J$-edge is never followed by an $L$-edge. Suppose two $L$-edges pass through $(a+i,m+i,1)$ for $i=0,1,2$. By the shape of each edge, $m+i\mid a+i$ and $R_a+4i\mid m+i+1$ for $i=0,1,2$. Put $K=a-m$. Then $m+i\mid K$ for each $i$, and one of $m,m+1,m+2$ is divisible by $3$, so $3\mid K$. Also $R_a+4i$ divides $4(m+i+1)-(R_a+4i)=4m+4-R_a=:C$, and one of $R_a,R_a+4,R_a+8$ is divisible by $3$, so $3\mid C$. But $C=4m+4-4a+p=p+4-4K\equiv p+1\pmod3$. Hence $p\equiv2\pmod3$. Two $J$-edges are excluded by the exchange.

(b) At $p=131$: $R_{33}=1$ divides $1+33$; $R_{34}=5$ divides $35$; $R_{35}=9$ divides $36$; each coordinate divides $pa$, and each pair is coprime. At $p=1613$ (prime): $3\mid405$, $7\mid406=7\cdot58$ and $11\mid407=11\cdot37$. Each step is $J(1,n)=(1,n+1)$.
\end{proof}
So the hypothesis $p\equiv1\pmod{12}$ of the supplement can be weakened to $p\not\equiv2\pmod3$. The next proposition shows that no weaker congruence hypothesis suffices, although the conclusion fails only at a minority of the primes $p\equiv2\pmod3$ (Remark~\ref{rem:sheardensity}). The primes $p\equiv2\pmod3$ are solvable in any case (Section~\ref{sec:core}); both propositions concern the transport graph, not solvability.

\begin{proposition}[the shear graph]\label{prop:sheargraph}
Let $p$ be a prime. Consider the directed graph whose vertices are the triples $(a,m,n)$ with $p/4<a<p$ and $(m,n)\in\mathcal W_a$, and whose arrows are the shear edges.
\begin{enumerate}[label=(\alph*)]
\item Every vertex has at most one outgoing and at most one incoming edge. So every connected component is a directed path.
\item \emph{Diagonal paths.} Let $R=4a-p$, and let $k\ge1$ with $a+k<p$. The vertices $(a+i,a+i,1)$, $0\le i\le k$, form a path of $L$-edges if and only if $R+4i$ divides $p+4$ for every $0\le i\le k$.
\item If $(a+i,m+i,1)$, $0\le i\le k$, is a path of $L$-edges and $R=4a-p$, then $R+4i$ divides $p+4$ for every $0\le i\le k-1$. In particular a path of two edges forces $R(R+4)\mid p+4$.
\item Every prime $p\equiv131\pmod{180}$ has a path of two edges. More generally, let $q\ge1$ and let $c$ be prime to $q$, with $c\equiv2\pmod3$ if $3\mid q$ (these are exactly the classes that contain primes $\equiv2\pmod3$). Then infinitely many primes $p\equiv c\pmod q$ have a path of two edges. So if Proposition~\ref{prop:twoshears}(a) holds at all but finitely many primes of a union of classes modulo $q$, then none of those classes contains primes $\equiv2\pmod3$.
\item For every $k\ge1$, infinitely many primes have a path of length $k$.
\end{enumerate}
\end{proposition}
\status{Proved here. Found by the referee pass of version~1; each part checked here. Part (a) was confirmed on the full graph of every prime below $1200$. A complete enumeration below $10^7$ through (c) and the shape of an edge finds:
\begin{itemize}
\item $19{,}558$ two-edge paths at $17{,}425$ primes;
\item $R(R+4)\mid p+4$ on every one of those paths;
\item all $13{,}822$ primes $p\equiv131\pmod{180}$ among the primes with a path.
\end{itemize}
The least primes with a path of length $1$, $2$, $3$, $4$ and $5$ are $11$, $131$, $3461$, $109{,}391$ and $1{,}183{,}451$. The referee gave $3{,}533{,}141$ for length $5$; it has such a path, but it is not the least.}
\begin{proof}
(a) By the shape of an edge, an $L$-edge starts at a vertex with $n=1$ and a $J$-edge at one with $m=1$. So the only possible successors of $(a,m,n)$ are $(a+1,m+1,1)$, when $n=1$, and $(a+1,1,n+1)$, when $m=1$. Both are possible only if $m=n=1$. Then $(1,1)\in\mathcal W_a$ forces $R_a\mid2$, so $R_a=1$. Either successor would then need $R_{a+1}=5$ to divide $3$, which is false.
The only possible predecessors of $(a+1,m',n')$ are $(a,m'-1,1)$, by an $L$-edge when $n'=1$, and $(a,1,n'-1)$, by a $J$-edge when $m'=1$. Both are possible only if $m'=n'=1$, and then $m'-1=0$ is not a positive coordinate.
Finally, $a$ increases along edges, so there are no cycles.

(b) The triple $(a+i,a+i,1)$ satisfies the divisibility and coprimality conditions of $\mathcal W_{a+i}$. It lies in $\mathcal W_{a+i}$ exactly when $R+4i\mid a+i+1$. As $R+4i$ is odd and $4(a+i+1)=(R+4i)+(p+4)$, this holds exactly when $R+4i\mid p+4$. Consecutive vertices are joined by an $L$-edge, since $L(a+i,1)=(a+i+1,1)$.

(c) Let $0\le i\le k-1$ and put $K=a-m$. The edge from the $i$-th vertex gives $m+i+1\mid a+i+1$, hence $m+i+1\mid K$. The $i$-th vertex lies in $\mathcal W_{a+i}$, so $R+4i\mid m+i+1$. Hence $R+4i\mid K$. Also $R+4i$ divides $4(m+i+1)-(R+4i)=4m+4-R$. Adding $4K$, it divides $4a+4-R=p+4$. For two edges, $R$ and $R+4$ are odd with difference $4$, hence coprime, so $R(R+4)\mid p+4$.

(d) If $p\equiv131\pmod{180}$, then $p\equiv3\pmod4$, so $R=1$ occurs, at $a=(p+1)/4$. Also $45\mid p+4$, so $1$, $5$ and $9$ divide $p+4$, and (b) applies with $k=2$.
In general, choose an odd $R$ with the following properties:
\begin{itemize}
\item $R\equiv-c\pmod4$ if $4\mid q$;
\item no prime $\ell\ge5$ dividing $q$ divides $R(R+4)(R+8)$;
\item the one of $R,R+4,R+8$ that is divisible by $3$ is not divisible by $9$.
\end{itemize}
Such an $R$ exists by the Chinese remainder theorem: each prime $\ell\ge5$ excludes three classes, and modulo $9$ the classes $0$, $1$ and $5$ are excluded.
Put $m'=R(R+4)(R+8)$. Its three factors are odd with pairwise differences $4$ and $8$, hence pairwise coprime, and $\gcd(m',q)\mid3$.
The congruences $p\equiv c\pmod q$, $p\equiv-4\pmod{m'}$ and $p\equiv-R\pmod4$ are compatible: modulo $3$ both of the first two give $p\equiv2$, and modulo $4$ the first and the third agree by the choice of $R$. Together they define a class prime to its modulus, which contains infinitely many primes by Dirichlet's theorem. Each such prime $p>R+8$ has the diagonal path of (b) with $a=(p+R)/4$ and $k=2$.
If $3\nmid q$, every class prime to $q$ contains primes $\equiv2\pmod3$ (by the Chinese remainder theorem and Dirichlet's theorem). If $3\mid q$, a class contains such primes exactly when $c\equiv2\pmod3$.

(e) Take $R=1$. The modulus $\operatorname{lcm}(1,5,9,\dots,4k+1)$ is odd, so the congruences $p\equiv3\pmod4$ and $p\equiv-4$ modulo it define a class prime to its modulus. Each prime $p>(4k+1)/3$ in it has, by (b), a diagonal path of length $k$ from $a=(p+1)/4$.
\end{proof}

\begin{remark}[how often the theorem fails]\label{rem:sheardensity}
Among the primes $p\equiv2\pmod3$, those with a path of two edges have upper relative density at most $0.335$. So the conclusion of Proposition~\ref{prop:twoshears}(a) holds on a set of primes $p\equiv2\pmod3$ of lower relative density at least $0.665$. Below $10^7$ it fails at $17{,}425$ of the $332{,}384$ primes $p\equiv2\pmod3$ ($5.24\%$); the smallest primes $p\equiv2\pmod3$ at which it holds are $2,5,11,17,23,29,41,47,53,59,71,83,\dots$.
\end{remark}
\begin{proof}
By Proposition~\ref{prop:twoshears}(a) and Proposition~\ref{prop:sheargraph}(c), a prime with a path of two edges satisfies:
\begin{itemize}
\item $p\equiv2\pmod3$;
\item for some odd $R$ with $R(R+4)\le p+4$, both $p\equiv-4\pmod{R(R+4)}$ and $p\equiv-R\pmod4$.
\end{itemize}
Fix $R$ and put $m=R(R+4)$. By Dirichlet's theorem, the second condition picks out a set of relative density $\delta_R=(1+[3\mid m])/(2\varphi(m))$ among the primes $p\equiv2\pmod3$.
For large $R$ the bound must be uniform in $x$:
\begin{itemize}
\item \emph{For $4m\le x^{1/2}$.} The Brun--Titchmarsh inequality $\pi(x;Q,b)\le2x/(\varphi(Q)\log(x/Q))$ for $Q<x$ \cite{MontgomeryVaughan}, with $Q=4m$, bounds the number of such primes up to $x$ by $2x/(\varphi(m)\log x)$. This is $(4+o(1))/\varphi(m)$ times the number of primes $p\equiv2\pmod3$ up to $x$.
\item \emph{For $4m>x^{1/2}$.} Each class has at most $x/m+1$ members up to $x$, and $R\le\sqrt{x+4}$. The total over these $R$ is $O(x^{3/4})$.
\end{itemize}
Since $\sum_R1/\varphi(R(R+4))$ converges, letting $x\to\infty$ and then the cut between small and large $R$ tend to infinity bounds the upper relative density by $\sum_{R\text{ odd}}\delta_R$.
Numerically, the partial sum over $R<2\cdot10^5$ is $0.334275$. The rest is less than $2\cdot10^{-4}$, by $n/\varphi(n)<e^\gamma\log\log n+3/\log\log n$ for $n\ge3$ \cite{RosserSchoenfeld}. So the sum is less than $0.335$.
\end{proof}

\subsection{The 124-page ES/Fable reader}\label{ssec:crosswalk}

It encodes a solution of $4/p=1/x+1/y+1/z$ as the binary quartic
\[
H_{\mathrm{ES}}(U,V)=-S^{-1}(U-pV)(U-xV)(U-yV)(U-zV),\qquad S=p+x+y+z,
\]
with normalised coefficients $u_0=-1/S$, $u_1=1$, $u_2=-\bigl(p(x+y+z)+xy+yz+zx\bigr)/S$, $u_3=\bigl(p(xy+yz+zx)+xyz\bigr)/S$ (which equals $5xyz/S$ on solutions) and $u_4=-pxyz/S$. It then transports the quartic to the Fable cover and to a weighted conductor of the zeta programme.

\begin{proposition}\label{prop:quartic}
\begin{enumerate}[label=(\alph*)]
\item On solutions, $p=-5u_4/u_3$ and
\[
625u_0u_4^3-125u_3u_4^2+25u_2u_3^2u_4-4u_3^4=0 .
\]
\item Conversely, let $u_0,u_2,u_3,u_4$ be complex numbers with $u_0u_3u_4\ne0$ that satisfy this relation. Put
\[
S=-1/u_0,\quad p=-5u_4/u_3,\quad P=u_3S/5,\quad \sigma_1=S-p,\quad \sigma_2=-u_2S-p\sigma_1 .
\]
Then the roots $x,y,z$ of $T^3-\sigma_1T^2+\sigma_2T-P$ are non-zero and satisfy $4/p=1/x+1/y+1/z$. The quartic $-S^{-1}(U-pV)(U-xV)(U-yV)(U-zV)$ has the coefficients $u_0,1,u_2,u_3,u_4$.
\end{enumerate}
\end{proposition}
\status{Proved here. The README of the package states (a), and \note{21\_} (Proposition~21.6) also checked the relation. Part (b) was found by the referee pass of version~1 and is proved here.}
\begin{proof}
(a) Write $\sigma_1=x+y+z$, $\sigma_2=xy+yz+zx$ and $P=xyz$. The equation is $4P=p\sigma_2$, so the third elementary symmetric function of $p,x,y,z$ is $p\sigma_2+P=5P$; this gives $u_3=5P/S$, and $p=-5u_4/u_3$ follows. Substituting and multiplying by $S^4/625$, the left side becomes $P^3\bigl(p^3-p^2S+p(p\sigma_1+\sigma_2)-4P\bigr)=P^3(p\sigma_2-4P)=0$.

(b) With the stated definitions, $u_0=-1/S$, $u_4=-pu_3/5$, $u_2=-(p\sigma_1+\sigma_2)/S$ and $u_3=5P/S$. Substituting, and using $S=p+\sigma_1$, gives
\[
625u_0u_4^3-125u_3u_4^2+25u_2u_3^2u_4-4u_3^4=\frac{u_3^3}{S}\bigl(5p^3-5p^2S+5p(p\sigma_1+\sigma_2)-20P\bigr)=\frac{5u_3^3}{S}\,(p\sigma_2-4P).
\]
So the relation gives $p\sigma_2=4P$. Now $xyz=P\ne0$, since $u_3\ne0$, and $p\ne0$, since $u_4\ne0$. Hence $1/x+1/y+1/z=\sigma_2/P=4/p$.
For the quartic, the coefficients of $U^4$, $U^3V$, $U^2V^2$, $UV^3$ and $V^4$ in $-S^{-1}(U-pV)(U-xV)(U-yV)(U-zV)$ are:
\begin{itemize}
\item $-1/S=u_0$;
\item $(p+\sigma_1)/S=1$;
\item $-(p\sigma_1+\sigma_2)/S=u_2$;
\item $(p\sigma_2+P)/S=5P/S=u_3$;
\item $-pP/S=u_4$.
\end{itemize}
\end{proof}
So the relation is the equation itself, written in these coordinates: by (a) and (b), it cuts out exactly the quartics that come from triples, possibly complex, with $4/p=1/x+1/y+1/z$. The Fable and conductor maps built on it are \textsf{located}; their comparison with the zeta programme is in the zeta audit (\note{47\_}, overlap OV-13).

\subsection{The continuation of 17--23 September}\label{ssec:continuation}

The README's \emph{Start reading} list has 28 items with proof packages in \texttt{research/incoming/}. The results bench R01--R10 and the integration arguments X1--X7 accompany them.
\begin{itemize}
\item The first-pass audit \note{43a\_} lists them all.
\item Section~\ref{sec:items} states the eleven items audited there.
\item Section~\ref{sec:negative} gives their negative results.
\end{itemize}

\subsection{The author's notes}\label{ssec:notescorpus}

The author's notes of April--June 2026 are separate from the workbench's documents, and Section~\ref{sec:notes} treats them. Their arithmetic is the shell calculus of Section~\ref{sec:core} together with the classical identity families. Their structural layers are Niemeier and umbral data, the class $6B$, Ogg's primes, Lorentzian geometry, Cayley--Dickson algebras and mock theta functions. The archive's §25 is drawn from one of them.

